\section{Introduction}
\label{sec:intro}

Quantum reservoir computing (QRC) is attractive because the quantum dynamics can remain fixed while only a classical readout is trained \cite{maass2002real,lukosevicius2009reservoir,tanaka2019recent,fujii2017harnessing,nakajima2019boosting,mujal2021opportunities}. In repeated-measurement QRC, however, the readout never sees an exact state. Each labeled time step is accessible only through a finite collection of repeated circuit executions. Once measurement cost becomes a dominant resource on near-term devices \cite{gonthier2022measurements}, the way this shot record is converted into classical training rows becomes part of the algorithm rather than a preprocessing detail.

The standard expectation-value (EV) protocol averages all shots for one time step into one clean feature vector. This is statistically attractive at the feature level, but it also compresses many circuit executions into only one readout-training row. Raw stacking is the opposite extreme: every shot becomes its own row, so the row count is maximal but each row carries the full measurement noise. In classical measurement-error terms, this is exactly the regime where noisy covariates attenuate fitted coefficients and flatten predictions \cite{fuller1987measurement,carroll2006measurement}. In finite-shot QRC this failure mode is visually obvious: forecasts become too damped even when the nominal row count is large.

This paper studies the middle ground, \emph{split-ensemble training}. The reservoir, observables, total execution budget, and final inference rule are all held fixed; only the organization of the shot record changes. At each labeled time step, the shots are partitioned into disjoint groups, averaged within each group, and used as several partially denoised training views of the same target. Compared with EV, this increases the number of readout-training rows. Compared with Raw, each row is substantially cleaner.

The main methodological question is \emph{when} grouping should help. We deliberately avoid a universal theorem for test error. Instead, we combine one exact protocol fact with a standard surrogate analysis used throughout classical measurement-error and high-dimensional ridge theory. The exact fact is simple: if EV has aspect ratio $\rhoEV=\Ttrain/d$, then grouping by $k$ shots yields $\rhoK=(\Nshots/k)\rhoEV$. Under a grouped-view surrogate, the grouped feature-noise covariance scales as $\Sigma_\epsilon/k$, while the effective row count scales as $\Nshots/k$. These ingredients produce an explicit attenuation-versus-row-count tradeoff and a cubic-root warm start for the preferred group size. The novelty is therefore not a new random-matrix theorem; it is the recognition that shot grouping gives an exact, budget-fair transformation of the readout geometry, and that standard theory is already sufficient to explain why the transformation matters.

The empirical story is intentionally organized around the strongest evidence rather than around a single hand-picked operating point. First, an empirical $\rhoEV$ sweep tests the regime picture directly. Second, stronger controls---EV-dup, diagonal variance-aware ridge baselines, and a larger-reservoir sweep---show that the effect is not exhausted by simple duplication or by a single small reservoir. Third, a hardware validation on IBM Quantum confirms that the same finite-shot protocol effect survives on a real backend.

A reviewer could reasonably worry that representative operating points with task-specific settings are too presentation-driven. We therefore make the representative comparison maximally uniform. Figure~\ref{fig:rep} and the baseline controls use one common synthetic operating point across all three tasks: same budget, same $N_{\mathrm{shots}}=18$, same $\sigma=0$, and same $\lambda=10$. This common point is chosen by a single explicit rule: it is the largest $N_{\mathrm{shots}}$ value in the shared sweep for which EV remains below $\NRMSE=1$ and Split still improves on EV on all three benchmarks in the $20$-seed mean. Figure~\ref{fig:rep} remains \emph{illustrative} rather than foundational, but it is now illustrative under one audited rule rather than under task-by-task exceptions.

The contributions are:
\begin{enumerate}[leftmargin=*,label=(\roman*)]
    \item We isolate \emph{shot organization under a fixed execution budget} as a distinct algorithmic axis in repeated-measurement QRC.
    \item We derive a grouped-view theory that combines the exact aspect-ratio lift $\rhoK=(\Nshots/k)\rhoEV$ with a standard attenuation-versus-row-count tradeoff and a closed-form warm start for the preferred group size.
    \item We validate the resulting regime picture on a corrected $20$-seed synthetic suite across Mackey--Glass, Lorenz, and NARMA10, emphasizing empirical gap curves and reviewer-facing controls rather than only representative points.
    \item We incorporate real-device validation on \texttt{ibm\_marrakesh} and show in additional simulations that the split advantage persists as the reservoir size increases through $Q=12$ for all three tasks, with positive paired deltas in every reported cell.
\end{enumerate}

The rest of the paper is organized as follows. Section~\ref{sec:background} defines the repeated-measurement setting, exact budget accounting, and the diagnostic $\rhoEV$. Section~\ref{sec:related} positions the protocol relative to finite-shot QRC and the classical measurement-error literature. Section~\ref{sec:methodology} presents the grouped-view protocol, the exact aspect-ratio lift, and the closed-form warm start. Section~\ref{sec:setup} describes the corrected synthetic protocol, the imported hardware experiment, and the larger-reservoir sweep. Section~\ref{sec:evaluation} presents the representative comparisons, empirical regime plot, baseline controls, hardware validation, and larger-reservoir results. Section~\ref{sec:conclusion} summarizes the practical decision rule and the main limitations.
